% Appendix: the EMBER instance (PLAN.md, decision 8 (v); docs/paper-mock.md, Section 4). The bookkeeping that Section 4.2
% summarizes: the item kinds, the readout forms with EMBER's features (from docs/evidence/code-audit-sound-directions.md,
% "Readout rule per EMBER feature"), the two features written over sums, and the re-encodings. Mostly the accepted text
% of the old Section 4.1 (session 27), and how the positions that locate items are recorded (block 3).

This appendix gives the details of the instance that \Cref{sec:instances:record} summarizes.

\paragraph{Items.} Scanned items are found by reading the bytes in order: single bytes, windows of 2,048 bytes taken
every 1,024 bytes for the entropy histogram, printable strings, and the Rich header, which the extractor searches for
within a fixed range of the file. Parsed items are found by following the file's structure: sections, the overlay
after the last section, header fields, imports, libraries, exports, the certificate table, and the entry section,
which contains the program's entry point.

\paragraph{Readouts.} \Cref{tab:ember:readouts} lists the features that EMBER reads through each form of
\Cref{tab:fw:readouts}. The monotone readouts are the presence of exports; the pools, which EMBER reads as 0 for a file
without sections; and single values that default to 0 when a file lacks the item, such as whether it has
relocations. EMBER computes two features in a way that the sum-then-readout form does not allow as written, and we
write them over sums without changing their values. EMBER divides each section's size by the file's size as it visits
the section, which would make the section's contribution depend on the rest of the file; we let the section
contribute its raw size and divide in the readout, where the file's size is a sum of its own. EMBER also keeps the
libraries as a set, so that a second import from a library already present adds nothing; we keep a count per library,
and the number of libraries becomes a distinct count.

\begin{table}[h]
  \caption{EMBER's features by the readout forms of \Cref{tab:fw:readouts}.}
  \label{tab:ember:readouts}
  \small
  \begin{tabular}{@{}lp{0.68\columnwidth}@{}}
    \toprule
    Readout & EMBER features \\
    \midrule
    identity & file size, counts of strings and sections, header fields, data directories, Rich header, overlay size,
      hashed buckets of sections, imports, and exports \\
    monotone & exports present, largest and smallest section entropy, has relocations, signature fields \\
    distinct count & number of libraries \\
    ratio & average string length, overlay and largest or smallest section size over file size \\
    normalized bin & byte histogram, byte-entropy histogram, printable characters \\
    any other & file and string entropy, machine type, subsystem, flag bits \\
    \bottomrule
  \end{tabular}
\end{table}

\paragraph{Re-encodings over sums.}
\label{sec:instances:sums}
Since the instrumented extractor records the sums, a detector can read them directly, each through an identity
readout, instead of the features. Many features move both ways only because of their readout: a normalized bin falls
when any other bin rises, while every count of the histogram can only rise when bytes are added. Reading the counts
keeps their directions. We use three such re-encodings: raw byte, string, and character counts in place of their
normalized frequencies; the same counts taken only over the section bodies and the overlay, where no in-place write of
our threat model lands; and the hashed section-name buckets, whose entries EMBER adds with a sign, split into their
positive and negative parts. Theorem~\ref{thm:fw:directions} gives their directions with no rule of their own, and the
same check covers them.

\paragraph{Positions that locate items.} For each kind, the record lists the positions the parser reads to find its
items: header fields and table entries, and positions that concern a whole kind, such as the number of sections and
the bounds of the range in which the extractor searches for the Rich header. Some are easy to miss. The overlay starts
where the last section or data directory ends, so the extent of every section and directory locates it. Others lie
past the end of the file, when the section table points to content that a truncated file no longer holds. We
completed the record by changing single bytes and checking whether the parse changed where no recorded position was
written. A first round found three parser rules the record had missed. After we recorded them, none of 58,666
single-byte changes across the headers and the directory data changed the parse without writing a recorded position.
